% Options for packages loaded elsewhere
\PassOptionsToPackage{unicode}{hyperref}
\PassOptionsToPackage{hyphens}{url}
\PassOptionsToPackage{dvipsnames,svgnames,x11names}{xcolor}
%
\documentclass[
  11pt,
  a4paper,
]{ltjsarticle}
\usepackage{amsmath,amssymb}
\usepackage{iftex}
\ifPDFTeX
  \usepackage[T1]{fontenc}
  \usepackage[utf8]{inputenc}
  \usepackage{textcomp} % provide euro and other symbols
\else % if luatex or xetex
  \usepackage{unicode-math} % this also loads fontspec
  \defaultfontfeatures{Scale=MatchLowercase}
  \defaultfontfeatures[\rmfamily]{Ligatures=TeX,Scale=1}
\fi
\usepackage{lmodern}
\ifPDFTeX\else
  % xetex/luatex font selection
\fi
% Use upquote if available, for straight quotes in verbatim environments
\IfFileExists{upquote.sty}{\usepackage{upquote}}{}
\IfFileExists{microtype.sty}{% use microtype if available
  \usepackage[]{microtype}
  \UseMicrotypeSet[protrusion]{basicmath} % disable protrusion for tt fonts
}{}
\makeatletter
\@ifundefined{KOMAClassName}{% if non-KOMA class
  \IfFileExists{parskip.sty}{%
    \usepackage{parskip}
  }{% else
    \setlength{\parindent}{0pt}
    \setlength{\parskip}{6pt plus 2pt minus 1pt}}
}{% if KOMA class
  \KOMAoptions{parskip=half}}
\makeatother
\usepackage{xcolor}
\usepackage[margin=24mm]{geometry}
\setlength{\emergencystretch}{3em} % prevent overfull lines
\providecommand{\tightlist}{%
  \setlength{\itemsep}{0pt}\setlength{\parskip}{0pt}}
\setcounter{secnumdepth}{5}
\ifLuaTeX
  \usepackage{selnolig}  % disable illegal ligatures
\fi
\IfFileExists{bookmark.sty}{\usepackage{bookmark}}{\usepackage{hyperref}}
\IfFileExists{xurl.sty}{\usepackage{xurl}}{} % add URL line breaks if available
\urlstyle{same}
\hypersetup{
  colorlinks=true,
  linkcolor={blue},
  filecolor={Maroon},
  citecolor={Blue},
  urlcolor={blue},
  pdfcreator={LaTeX via pandoc}}

\author{}
\date{}

\begin{document}

\hypertarget{span}{%
\section{Span}\label{span}}

\texttt{v\_1,...,v\_k} のすべての一次結合の集合を

\[
\operatorname{span}\{v_1,\dots,v_k\}
\]

と書きます。

\hypertarget{ux5e7eux4f55ux5b66ux7684ux610fux5473}{%
\subsection{幾何学的意味}\label{ux5e7eux4f55ux5b66ux7684ux610fux5473}}

\texttt{R\^{}3} で

\begin{itemize}
\tightlist
\item
  1本の非零ベクトルの span は原点を通る直線
\item
  独立な2本の span は原点を通る平面
\item
  独立な3本の span は \texttt{R\^{}3} 全体
\end{itemize}

になります。

ここで重要なのは、ベクトルの「本数」ではなく、独立な方向の数です。同じ方向のベクトルを何本追加しても
span は広がりません。

\hypertarget{ux9023ux7acbux65b9ux7a0bux5f0fux3068ux306eux63a5ux7d9a}{%
\subsection{連立方程式との接続}\label{ux9023ux7acbux65b9ux7a0bux5f0fux3068ux306eux63a5ux7d9a}}

\texttt{Ax=b} が解を持つ条件は

\[
b\in\operatorname{span}\{a_1,\dots,a_n\}
\]

です。つまり \texttt{b} が \texttt{A} の列空間に存在することです。

この一文だけで、連立方程式の可解性を幾何学的に表現できます。

\end{document}
